\chapter{Fundamentação Teórica}
\label{cap:fundamentacao}

Este capítulo apresenta os principais conceitos necessários para a compreensão deste trabalho. Inicialmente, são abordados aspectos relacionados à ovinocultura e ao monitoramento de animais em sistemas de produção. Em seguida, são apresentados conceitos relacionados à pecuária de precisão, à Internet das Coisas e ao uso de sensores no contexto agropecuário. Por fim, são discutidas redes de comunicação de longo alcance e baixo consumo, com ênfase na tecnologia \textit{Long Range Wide Area Network} (LoRaWAN), na escalabilidade da comunicação e na simulação de redes aplicadas ao monitoramento de ovinos.

\section{Ovinocultura}
\label{sec:ovinocultura}

A ovinocultura é uma atividade voltada à criação de ovinos, podendo ser direcionada à produção de carne, leite e outros produtos de origem animal. No contexto do Nordeste e do Ceará, essa atividade apresenta importância econômica, social e cultural, especialmente em regiões sertanejas. O perfil da cadeia da ovinocaprinocultura cearense aponta que o Ceará possui relevância nesse setor, sendo o quarto maior rebanho de caprinos e ovinos do Brasil, além de apresentar potencial de crescimento, sobretudo no sertão \cite{ceara2025}.

No Semiárido brasileiro, a criação de caprinos e ovinos está associada à organização produtiva regional, contribuindo para a geração de ocupação, produção de alimento e complementação da renda de famílias rurais. Além disso, a atividade ocorre em um ambiente marcado por limitações climáticas, variações na disponibilidade de água e necessidade de adaptação dos sistemas produtivos às condições locais \cite{voltolini2011}. Essas características ajudam a compreender a diversidade dos ambientes de criação e servem de base para a definição de cenários distintos de manejo na simulação proposta neste trabalho.

A criação de ovinos na região semiárida também pode estar associada a sistemas tradicionais de produção, nos quais os animais permanecem em áreas de pastagem durante parte do dia e podem ser recolhidos em instalações, como apriscos, em determinados períodos. No estudo de \citeonline{brandao2024}, por exemplo, os cordeiros permaneceram com suas mães em piquetes durante o dia e no aprisco durante a noite, caracterizando um manejo comum na região. Esse tipo de organização mostra que os animais podem alternar entre momentos de maior concentração e momentos de maior deslocamento, aspecto importante para este trabalho, pois os cenários de confinamento e pastagem serão considerados na simulação da rede de monitoramento.

Além disso, a produção ovina no Semiárido pode ser influenciada por fatores como disponibilidade de alimento, período seco, condições ambientais e forma de manejo adotada. Em sistemas baseados em pastagens nativas, a variação na quantidade e na qualidade da forragem pode repercutir no desempenho produtivo dos animais, enquanto o confinamento e o uso de dietas adaptadas ao Semiárido, como aquelas à base de palma forrageira, podem contribuir para melhorar a produtividade \cite{cartaxo2025}. Nesse contexto, o acompanhamento de variáveis relacionadas ao rebanho e ao ambiente de criação torna-se relevante, pois estudos em pecuária de precisão com ovinos indicam que dados ambientais, fisiológicos e comportamentais podem apoiar a avaliação do bem-estar, da produtividade e da tomada de decisão no manejo \cite{nikolopoulou2026}. Dessa forma, compreender as características básicas da ovinocultura no Semiárido é importante para representar, de maneira coerente, os cenários utilizados neste trabalho.

\section{Pecuária de Precisão e Internet das Coisas}
\label{sec:pecuaria-precisao-iot}

A pecuária de precisão consiste no uso de tecnologias para acompanhar, de forma mais contínua e objetiva, informações relacionadas aos animais, ao ambiente de criação e ao manejo produtivo. Essa abordagem busca apoiar o produtor na tomada de decisão, permitindo que alterações no comportamento, na saúde ou nas condições do ambiente sejam identificadas com maior rapidez. De acordo com \citeonline{terence2024}, a pecuária de precisão, do inglês \textit{Precision Livestock Farming} (PLF), utiliza tecnologias para avaliar animais, comportamento, produção e ambiente, podendo fornecer informações ao produtor ou acionar sistemas automáticos de controle. Dessa forma, a pecuária de precisão contribui para tornar o manejo mais eficiente, principalmente em sistemas nos quais o acompanhamento manual dos animais pode ser limitado pela quantidade de indivíduos, pela distância ou pela disponibilidade de mão de obra.

A Internet das Coisas, do inglês \textit{Internet of Things} (IoT), permite a conexão de sensores, dispositivos e sistemas capazes de coletar e transmitir dados sobre o ambiente de produção. Na agricultura de precisão, essa integração entre sensores inteligentes, IoT e análise de dados contribui para o acompanhamento de variáveis importantes e para o uso mais eficiente dos recursos disponíveis. \citeonline{soussi2024} destacam que a agricultura de precisão utiliza sensores inteligentes e tecnologias como IoT, análise de grandes volumes de dados e inteligência artificial para orientar decisões e otimizar práticas agrícolas. De forma complementar, \citeonline{rajak2023} apontam que sensores associados à IoT podem monitorar variáveis como umidade, temperatura, composição do solo e outros fatores ambientais relevantes para a produção. No contexto da pecuária, essa mesma lógica pode ser aplicada ao acompanhamento de animais e ambientes de criação, permitindo que dados coletados em campo apoiem decisões relacionadas ao manejo e ao monitoramento do rebanho.

No monitoramento animal, sensores podem ser utilizados para coletar informações fisiológicas, comportamentais e ambientais, permitindo acompanhar o estado dos animais de forma mais contínua. Essas informações podem incluir temperatura corporal, frequência cardíaca, movimentação, localização e condições do ambiente de criação. \citeonline{nikolopoulou2026} destacam que, em sistemas de pecuária de precisão voltados a ovinos, dados provenientes de sensores ambientais, fisiológicos e comportamentais podem apoiar a identificação de alterações relacionadas ao bem-estar e à produtividade dos animais. De forma semelhante, \citeonline{bhaskaran2024} apresentam o uso de sensores e sistemas de IoT para acompanhar indicadores como padrões de movimento, temperatura corporal e frequência cardíaca em sistemas de monitoramento da saúde animal. Dessa forma, o uso de sensores na pecuária permite transformar dados coletados em campo em informações úteis para o acompanhamento do rebanho e para o apoio à tomada de decisão.

\section{Redes LPWAN, LoRa e LoRaWAN}
\label{sec:lpwan-lora-lorawan}

As aplicações de Internet das Coisas em ambientes rurais dependem de tecnologias de comunicação capazes de operar em áreas extensas, com baixo consumo de energia e custo reduzido. Nesse contexto, as redes do tipo \textit{Low-Power Wide-Area Network} (LPWAN) são utilizadas em aplicações que exigem comunicação de longo alcance para dispositivos com restrições de energia, como sensores alimentados por bateria. \citeonline{pagano2023} destacam que uma das dificuldades para a digitalização da agricultura está relacionada à limitação de conectividade em muitas áreas rurais, sendo as LPWANs uma alternativa adequada por apresentarem características como ampla cobertura, baixo consumo de energia e baixo custo operacional. De forma semelhante, \citeonline{jouhari2023} apontam que as LPWANs se consolidaram como uma opção de conectividade para aplicações de IoT, especialmente por permitirem a comunicação de dispositivos de baixa potência em cenários nos quais a transmissão de pequenos volumes de dados é suficiente.

A tecnologia \textit{Long Range} (LoRa) corresponde à camada física de comunicação, sendo responsável pela transmissão sem fio de longo alcance e baixo consumo de energia. Já a \textit{Long Range Wide Area Network} (LoRaWAN) define a camada de acesso ao meio e a arquitetura de rede utilizada para organizar a comunicação entre os dispositivos finais, os \textit{gateways} e os servidores de rede. Segundo \citeonline{jouhari2023}, LoRa é considerada a camada física, enquanto LoRaWAN corresponde à camada de controle de acesso ao meio, adotando uma topologia em estrela para permitir a comunicação entre múltiplos dispositivos finais e o \textit{gateway}. Essa organização é adequada para aplicações de IoT em que sensores precisam transmitir pequenas quantidades de dados periodicamente, como ocorre em sistemas de monitoramento ambiental, agrícola e pecuário.

A arquitetura da LoRaWAN é formada por dispositivos finais, \textit{gateways}, servidor de rede e servidores de aplicação. Os dispositivos finais, também chamados de nós sensores, são responsáveis por coletar dados e transmiti-los por comunicação sem fio. Os \textit{gateways} recebem os pacotes enviados pelos dispositivos finais e os encaminham para o servidor de rede, que realiza o gerenciamento da comunicação e o tratamento das mensagens recebidas. Em seguida, os dados podem ser disponibilizados aos servidores de aplicação, nos quais as informações são armazenadas, processadas ou apresentadas ao usuário. De acordo com \citeonline{jouhari2023}, a LoRaWAN utiliza uma topologia do tipo estrela de estrelas, na qual os dispositivos finais se comunicam com os \textit{gateways}, e estes encaminham os dados para o servidor de rede. Essa estrutura é importante para aplicações de monitoramento, pois permite que vários sensores distribuídos em uma área enviem informações para uma infraestrutura centralizada.

A LoRaWAN define diferentes classes de dispositivos finais, sendo elas Classe A, Classe B e Classe C. A Classe A é a classe básica e deve ser implementada por todos os dispositivos finais, permitindo comunicação bidirecional com baixo consumo de energia. Nesse modo de operação, cada transmissão de subida, conhecida como \textit{uplink}, é seguida por duas janelas curtas de recepção, nas quais o dispositivo pode receber mensagens de descida, chamadas de \textit{downlink}. Como o dispositivo permanece a maior parte do tempo em estado de repouso e só abre janelas de recepção após transmitir dados, essa classe apresenta menor consumo de energia, sendo adequada para aplicações com sensores alimentados por bateria \cite{loraalliance2020}. As Classes B e C, por sua vez, ampliam as possibilidades de recepção de mensagens, mas aumentam o consumo energético, pois exigem janelas de recepção adicionais ou praticamente contínuas.

Em redes LoRaWAN, o desempenho da comunicação depende da configuração de diferentes parâmetros de transmissão. Entre eles, destacam-se o fator de espalhamento, do inglês \textit{Spreading Factor} (SF), a potência de transmissão, do inglês \textit{Transmission Power} (TP), a largura de banda, do inglês \textit{Bandwidth} (BW), e a taxa de codificação, do inglês \textit{Coding Rate} (CR). O SF influencia diretamente o alcance e a taxa de dados: valores mais altos permitem maior alcance, porém reduzem a taxa de transmissão e aumentam o tempo de ocupação do canal. A potência de transmissão interfere na intensidade do sinal enviado, podendo ampliar o alcance da comunicação, mas também aumentar o consumo energético. A largura de banda afeta a velocidade de transmissão e a sensibilidade do sinal, enquanto a taxa de codificação adiciona redundância aos dados transmitidos, contribuindo para maior confiabilidade em ambientes com ruído, embora reduza a taxa efetiva de transmissão. Assim, a configuração desses parâmetros deve buscar equilíbrio entre alcance, consumo energético, confiabilidade e capacidade da rede \cite{sarmentoneto2025}.

O mecanismo de \textit{Adaptive Data Rate} (ADR) é utilizado em redes LoRaWAN para ajustar dinamicamente parâmetros de transmissão dos dispositivos finais, como o fator de espalhamento e a potência de transmissão. Esse ajuste busca melhorar o desempenho da comunicação, reduzir o consumo energético e utilizar os recursos da rede de forma mais eficiente. De acordo com \citeonline{sarmentoneto2025}, o ADR estima a qualidade do canal de comunicação a partir de valores como a relação sinal-ruído, do inglês \textit{Signal-to-Noise Ratio} (SNR), e, com base nessa estimativa, adapta os parâmetros de transmissão do dispositivo. Dessa forma, dispositivos com boas condições de comunicação podem utilizar configurações mais eficientes, enquanto dispositivos em condições menos favoráveis podem necessitar de parâmetros que aumentem o alcance da transmissão.

Apesar de sua importância para o ajuste automático da comunicação, o ADR apresenta limitações quando aplicado a dispositivos móveis. Isso ocorre porque o mecanismo foi projetado principalmente para dispositivos estáticos, nos quais as condições do canal tendem a variar de forma mais lenta. Em cenários móveis, como no monitoramento de animais em pastagem, a posição dos nós sensores pode mudar com frequência, alterando a distância em relação ao \textit{gateway} e as condições de propagação do sinal. Segundo \citeonline{sarmentoneto2025}, em ambientes móveis, as variações frequentes do canal podem reduzir o desempenho do ADR padrão, aumentando a perda de pacotes e o consumo energético quando a adaptação dos parâmetros não acompanha adequadamente essas mudanças. Dessa forma, a mobilidade dos dispositivos deve ser considerada em estudos que avaliam o desempenho e a escalabilidade de redes LoRaWAN aplicadas ao monitoramento animal.

\section{Escalabilidade em redes LoRaWAN}
\label{sec:escalabilidade-lorawan}

A escalabilidade de uma rede está relacionada à sua capacidade de manter um desempenho adequado mesmo quando ocorre aumento na quantidade de dispositivos conectados, no volume de transmissões ou na área de cobertura. Em redes LoRaWAN, esse aspecto é importante porque diversos dispositivos finais podem compartilhar o mesmo meio de comunicação para transmitir pequenos pacotes de dados ao \textit{gateway}. \citeonline{jouhari2023} destacam que, apesar da ampla utilização da tecnologia LoRa em aplicações de IoT, a interferência entre sinais LoRa e as colisões decorrentes de transmissões simultâneas estão entre as principais limitações para a ampliação da capacidade da rede. Dessa forma, em cenários com muitos nós sensores transmitindo periodicamente, como em aplicações de monitoramento animal, o aumento da quantidade de dispositivos pode afetar diretamente a entrega dos pacotes e a eficiência da comunicação.

Em aplicações de monitoramento animal, a escalabilidade da rede pode ser afetada por fatores como a quantidade de nós sensores, o intervalo de envio das mensagens, a distância entre os dispositivos e o \textit{gateway}, a mobilidade dos animais e a configuração dos parâmetros de transmissão. Quando muitos dispositivos transmitem periodicamente, aumenta a disputa pelo meio de comunicação, podendo ocorrer colisões, perdas de pacotes e redução da eficiência da rede. \citeonline{sarmentoneto2025} destacam que, em redes LoRaWAN, o aumento da densidade de dispositivos e a mobilidade dos nós são fatores que influenciam diretamente métricas como taxa de entrega de pacotes, eficiência energética, vazão e colisões. Assim, em um cenário no qual cada ovino é representado por um nó sensor, a análise da escalabilidade permite observar até que ponto a rede consegue manter a comunicação adequada conforme o número de animais monitorados aumenta.

A avaliação da escalabilidade em redes LoRaWAN pode ser realizada por meio de métricas que indicam a qualidade da comunicação e o comportamento dos dispositivos durante a transmissão dos dados. Entre essas métricas, destaca-se a taxa de entrega de pacotes, do inglês \textit{Packet Delivery Ratio} (PDR), que representa a proporção de pacotes enviados pelos dispositivos finais que são recebidos com sucesso pelo \textit{gateway} ou pelo servidor da rede. Além da PDR, também podem ser observadas perdas de pacotes, colisões, vazão da rede, do inglês \textit{throughput}, e consumo energético dos dispositivos. Em estudos de simulação, essas métricas permitem verificar como a rede se comporta em diferentes condições, como aumento da quantidade de nós sensores, variação na posição dos dispositivos e mudanças nos parâmetros de transmissão. \citeonline{macedo2023} utilizam métricas como PDR, perda de sinal e \textit{throughput} para avaliar o desempenho de uma rede LoRaWAN aplicada ao monitoramento pecuário, enquanto \citeonline{sarmentoneto2025} consideram PDR, eficiência energética, vazão e colisões na análise de escalabilidade de redes LoRaWAN móveis.

\section{Simulação de redes LoRaWAN}
\label{sec:simulacao-lorawan}

A simulação de redes é uma abordagem utilizada para avaliar o comportamento de sistemas de comunicação antes de sua implantação em ambiente real. Em redes LoRaWAN, a simulação permite representar diferentes quantidades de dispositivos, áreas de cobertura, padrões de mobilidade, intervalos de transmissão e parâmetros de comunicação, possibilitando observar como esses fatores influenciam o desempenho da rede. Essa abordagem é especialmente útil em cenários de monitoramento animal, pois a realização de testes práticos com grande quantidade de dispositivos pode exigir maior custo, tempo e infraestrutura. Dessa forma, a simulação computacional permite analisar previamente o funcionamento da rede e identificar possíveis limitações relacionadas à escalabilidade, à entrega de pacotes e ao consumo energético.

Em uma simulação de rede LoRaWAN, é necessário definir os principais elementos que compõem o cenário analisado, como os dispositivos finais, os \textit{gateways}, a área de cobertura, o intervalo de transmissão das mensagens, os parâmetros de comunicação e o comportamento dos nós ao longo do tempo. No caso de aplicações voltadas ao monitoramento animal, os dispositivos finais podem representar sensores instalados nos animais, responsáveis por enviar periodicamente informações para a infraestrutura da rede. \citeonline{macedo2023}, por exemplo, modelam os animais como dispositivos finais em uma rede LoRaWAN, considerando diferenças entre cenários de confinamento e pastagem. Já \citeonline{sarmentoneto2025} utilizam simulações para avaliar redes LoRaWAN com diferentes densidades de dispositivos e modelos de mobilidade, permitindo observar o impacto dessas variáveis sobre o desempenho da comunicação. Dessa forma, a simulação possibilita representar cenários próximos ao contexto de aplicação desejado e avaliar como alterações no número de nós, na mobilidade e na configuração da rede influenciam os resultados obtidos.

Diante dos conceitos apresentados, observa-se que o monitoramento de ovinos por meio de sensores depende não apenas da coleta dos dados, mas também da capacidade da rede de comunicação em transmitir essas informações de forma eficiente. A ovinocultura no Semiárido apresenta características de manejo que podem envolver tanto a concentração dos animais em espaços reduzidos quanto o deslocamento em áreas de pastagem. Nesse contexto, tecnologias como LoRaWAN tornam-se relevantes por permitirem comunicação de longo alcance e baixo consumo energético, características adequadas para aplicações rurais. Entretanto, fatores como quantidade de nós sensores, mobilidade, intervalo de transmissão, colisões e eficiência energética podem influenciar o desempenho da rede. Assim, os conceitos de pecuária de precisão, Internet das Coisas, LoRaWAN, escalabilidade e simulação de redes formam a base teórica necessária para a análise proposta neste trabalho.